%==============================================================
% Making AI More Trustworthy: Context Labelling and
% Original-Source Tracking for Clinical Question Answering
%
% Thesis-format report. Compile on Overleaf (recommended) or
% locally with:
%   pdflatex main  ->  bibtex main  ->  pdflatex main  x2
% Document class: report (available by default on Overleaf).
%==============================================================
\documentclass[11pt,a4paper]{report}

\usepackage[utf8]{inputenc}
\usepackage[T1]{fontenc}
\usepackage[margin=1in]{geometry}
\usepackage{cite}
\usepackage{amsmath,amssymb,amsfonts}
\usepackage{graphicx}
\usepackage{textcomp}
\usepackage{xcolor}
\usepackage{booktabs}
\usepackage{multirow}
\usepackage{array}
\usepackage{url}
\usepackage{caption}
\usepackage{float}
\usepackage[hidelinks]{hyperref}

% Look in the local figures/ folder first (self-contained for Overleaf),
% then fall back to the repository outputs folder.
\graphicspath{{figures/}{./figures/}{../outputs/figures/paper/}}

% Concise, brand-neutral references to the system.
\newcommand{\zt}{\mathbf{z}_t}

\begin{document}

%==============================================================
% TITLE PAGE
%==============================================================
\begin{titlepage}
\centering
\vspace*{1.5cm}
{\LARGE\bfseries Making AI More Trustworthy:\\[0.35em]
Context Labelling and Original-Source Tracking\\[0.35em]
for Clinical Question Answering\par}
\vspace{1.4cm}
{\large A World-Model--Guided, Calibrated, and Interpretable\\
Retrieval-Augmented Generation Framework\par}
\vspace{2.2cm}
{\large\itshape Author Name\par}
\vspace{0.4cm}
{\normalsize Department / Institution\par}
{\normalsize City, Country\par}
\vspace{2.4cm}
{\normalsize A thesis submitted in partial fulfilment\\
of the requirements for the degree of\\
[Degree Title]\par}
\vfill
{\normalsize \today\par}
\end{titlepage}

%==============================================================
% ABSTRACT
%==============================================================
\chapter*{Abstract}
\addcontentsline{toc}{chapter}{Abstract}

Large language models (LLMs) encode substantial clinical knowledge, yet they
remain difficult to trust: they hallucinate fluent but unsupported claims, they
rarely expose a calibrated sense of their own uncertainty, and they seldom make
clear \emph{which} pieces of evidence drove an answer. Retrieval-augmented
generation (RAG) reduces hallucination by grounding generation in retrieved
documents, but the canonical \emph{retrieve--concatenate--generate} pipeline is
reactive: it neither models the consequences of clinical decisions nor tracks
the provenance of each statement. This thesis presents a framework that makes
clinical question answering more trustworthy through two coupled mechanisms:
\textbf{context labelling}---attaching structured, calibrated labels (retrieved
evidence, causal knowledge-graph relations, a predicted diagnostic state, and an
uncertainty estimate) to every answer---and \textbf{original-source
tracking}---tracing each claim back to the specific source documents that
supported it using Shapley-value attribution and explicit citations.

The framework augments a mature hybrid retrieval stack with three trust-oriented
capabilities: (i) an action-conditioned \emph{latent world model} that simulates
how a patient state evolves and supports deliberate, tree-of-thought plan
search; (ii) a deep-ensemble predictor that produces \emph{calibrated}
diagnostic uncertainty with an explicit epistemic/aleatoric decomposition; and
(iii) a Shapley-value attribution module that quantifies the marginal
contribution of every retrieved document. The neural core---a causal
graph-attention encoder, a gated recurrent latent-dynamics model with a
supervised reward head, and a five-member probabilistic ensemble---contains only
$10.1$M parameters and trains in under $20$ minutes on a single consumer GPU. On
a corpus assembled from DDXPlus, MedMCQA, and MedQA ($100{,}000$ samples), the
world model attains a next-state mean absolute error of $0.039$ and an expected
calibration error of $0.029$, with a train-to-test error gap below $0.002$ MAE,
indicating strong generalisation. In a controlled comparison against a strong
hybrid (dense $+$ BM25) RAG baseline on ten complex clinical vignettes, the
proposed framework improves clinical concept coverage from $0.70$ to $0.97$,
answer-structure completeness from $0.50$ to $0.84$, and evidence grounding from
$2.0$ to $5.1$ citations per answer; the improvement is statistically
significant (Wilcoxon signed-rank $p=0.046$, mean $+0.27$, $95\%$ CI $\pm0.21$).
Against four advanced-RAG pipelines (vanilla, hybrid, HyDE, and
cross-encoder$+$MMR) run live through the same LLM, the framework leads on every
answer-quality metric while uniquely providing calibrated uncertainty,
per-document source attribution, and treatment-plan simulation.
All code, configurations, and figure-generating scripts are released to support
reproducibility.

\vspace{0.5em}
\noindent\textbf{Keywords:} trustworthy AI, retrieval-augmented generation,
world models, uncertainty calibration, Shapley attribution, source tracking,
clinical question answering.

\tableofcontents

%==============================================================
\chapter{Introduction}
%==============================================================

\section{Background}
Clinical decision support is among the most consequential applications of
artificial intelligence, and large language models (LLMs) have made fluent
medical question answering appear within reach: they encode a remarkable amount
of textbook and guideline knowledge~\cite{singhal2023medpalm}. Yet fluency is
not trustworthiness. Two failure modes block safe deployment. First, LLMs
\emph{hallucinate}---they produce confident, well-formed claims that are not
supported by any source~\cite{ji2023hallucination}. Second, they are
\emph{opaque about confidence and provenance}: a clinician cannot easily tell
whether an answer rests on solid evidence or on a plausible-sounding
fabrication, nor which documents were actually decisive.

Retrieval-augmented generation (RAG)~\cite{lewis2020rag} addresses the first
failure mode by conditioning generation on retrieved documents and is now
standard practice. Modern systems combine sparse lexical scoring
(BM25~\cite{robertson2009bm25}) with dense
bi-encoders~\cite{reimers2019sbert}, fuse the rankings~\cite{cormack2009rrf},
and re-rank with cross-encoders~\cite{nogueira2019passage}. However, the
canonical \emph{retrieve--concatenate--generate} pipeline remains
\emph{reactive}: it treats retrieved passages as a static context window and
delegates all reasoning, self-assessment, and provenance to the opaque forward
pass of the LLM. For high-stakes clinical use this is not enough.

\section{Rationale}
Trust in a clinical AI system rests on three properties that the canonical RAG
pipeline does not provide. \emph{Prospective reasoning:} good clinical practice
simulates how a patient will respond to candidate actions before committing to
one. \emph{Calibrated self-knowledge:} a trustworthy system must know what it
does not know, so that an answer given with $0.95$ confidence and one given at
the edge of the training distribution do not look identical. \emph{Provenance:}
for auditing and accountability, a system must state \emph{which} evidence was
decisive, not merely cite documents decoratively.

This thesis argues that these three properties can be delivered by reframing the
problem around two mechanisms. \textbf{Context labelling} attaches a structured,
machine-checkable record to every answer: the retrieved evidence set, the
causal knowledge-graph relations consulted, the predicted diagnostic state, and
a calibrated uncertainty estimate. \textbf{Original-source tracking} then traces
each claim to the specific documents that supported it, using Shapley-value
attribution to quantify each document's marginal contribution. Together, these
mechanisms convert an opaque generator into an auditable, calibrated, and
source-grounded assistant.

\section{Problem Statement}
Given a clinical query $q$ and a corpus of medical documents, the goal is to
produce an answer that is (a) grounded in retrieved evidence, (b) accompanied by
a calibrated confidence and an explicit uncertainty decomposition, (c)
attributable to the specific source documents that determined it, and (d)
informed by a forward simulation of the patient's diagnostic trajectory rather
than a single reactive pass. Formally, we require a system that jointly outputs
an answer, a set of context labels (evidence, graph relations, predicted state,
uncertainty), and a per-document attribution vector, while remaining compact
enough to train and audit on commodity hardware.

\section{Objectives}
The objectives of this work are to:
\begin{enumerate}
\item design a RAG framework that augments hybrid retrieval with an
action-conditioned latent world model for prospective, tree-of-thought planning;
\item produce \emph{calibrated} predictive uncertainty with an explicit
epistemic/aleatoric decomposition via a deep ensemble;
\item implement \emph{original-source tracking} through Shapley-value evidence
attribution and explicit citations;
\item keep the neural core compact ($\sim$$10$M parameters) and trainable in
under twenty minutes on a single consumer GPU;
\item evaluate trustworthiness with calibration, generalisation, statistical
significance, and a controlled comparison against a strong RAG baseline; and
\item characterise the method's scope and limitations honestly to support
responsible interpretation.
\end{enumerate}

\section{Methodology in Brief}
The proposed framework is organised as a twelve-stage pipeline grouped into five
functional blocks: \emph{retrieval}, \emph{graph reasoning}, \emph{world-model
planning}, \emph{uncertainty}, and \emph{explanation}
(Fig.~\ref{fig:pipeline}). A query is expanded and answered by hybrid retrieval
and re-ranking; the grounded evidence and a causal knowledge graph are encoded
into a latent patient state; the latent world model rolls the state forward
under a tree-of-thought beam planner; a deep ensemble scores the outcome with
calibrated uncertainty; a hallucination check verifies the draft; and a
Shapley-value module attributes the answer to individual documents before the
LLM verbalises the final, source-tracked response. Figure~\ref{fig:archdetail}
gives the detailed neural architecture with tensor shapes and shows where
context labels enter and where source tracking is produced.

\begin{figure}[H]
  \centering
  \includegraphics[width=\textwidth]{fig1_pipeline}
  \caption{Inference pipeline for trustworthy clinical question answering. A
  query is expanded, answered by hybrid retrieval and re-ranking, encoded with a
  causal knowledge-graph GNN into a latent state, rolled forward by the latent
  world model under a tree-of-thought beam planner, scored for calibrated
  uncertainty by a deep ensemble, verified for hallucination, attributed to
  evidence via Shapley values, and finally verbalised by the LLM with explicit
  source tracking.}
  \label{fig:pipeline}
\end{figure}

\begin{figure}[H]
  \centering
  \includegraphics[width=\textwidth]{fig17_architecture_detailed}
  \caption{Detailed neural architecture with tensor shapes. \emph{Context
  labels} (patient observations, action labels, and knowledge-graph relations)
  enter on the left; the causal GNN, evidence-fusion encoder, action-conditioned
  causal residual, and recurrent latent dynamics form the $10.1$M-parameter
  core; and the right-hand heads, beam planner, deep ensemble, and Shapley
  \emph{source-tracking} module produce a calibrated, source-attributed answer.}
  \label{fig:archdetail}
\end{figure}

\section{Scope and Challenges}
The scope of this thesis is the design, implementation, and evaluation of the
trust mechanisms (context labelling and source tracking) and the world-model
core that supports them, evaluated on public clinical-QA resources. Several
challenges shape the work. The dynamics are trained on \emph{semi-synthetic}
trajectories constructed by revealing evidence sequentially, so the transition
target is close to deterministic. The actions correspond to evidence
acquisition rather than physical interventions. The knowledge graph encodes a
stochastic causal \emph{prior} rather than a curated ontology. Fine-grained,
imbalanced rank classification over $49$ pathologies makes macro-F1 a hard
metric. We treat each of these as an explicit boundary of validity rather than a
hidden assumption, and return to them in Sections~\ref{sec:discussion}
and~\ref{sec:future}.

%==============================================================
\chapter{Literature Review}
%==============================================================

\section{Retrieval-Augmented Generation}
RAG~\cite{lewis2020rag} and dense passage retrieval~\cite{karpukhin2020dpr}
established the practice of conditioning generation on retrieved evidence.
Production systems combine sparse lexical scoring
(BM25~\cite{robertson2009bm25}) with dense bi-encoders~\cite{reimers2019sbert}
indexed for fast search~\cite{johnson2019faiss}, fuse the ranked lists with
reciprocal rank fusion~\cite{cormack2009rrf}, and re-rank with
cross-encoders~\cite{nogueira2019passage} while promoting diversity through
maximal marginal relevance~\cite{carbonell1998mmr}. Hypothetical document
embeddings (HyDE)~\cite{gao2023hyde} improve recall by expanding queries. Our
framework adopts this mature retrieval stack but does not stop there: it feeds
the grounded state into a world model and tracks the provenance of every
retrieved document.

\section{World Models and Planning}
World models learn the dynamics of an environment in a latent space and plan by
imagining rollouts~\cite{ha2018worldmodels}. The Dreamer
family~\cite{hafner2020dreamer,hafner2023dreamerv3} shows that latent recurrent
dynamics support sample-efficient control. Tree-of-thought
reasoning~\cite{yao2023tot} and chain-of-thought prompting~\cite{wei2022cot}
bring deliberate, multi-step search to LLMs. We connect these threads by using a
learned clinical latent-dynamics model as the transition function for a
tree-of-thought beam search, turning reactive answering into prospective
reasoning.

\section{Uncertainty and Interpretability}
Deep ensembles~\cite{lakshminarayanan2017ensembles} remain a strong, simple
baseline for predictive uncertainty, and modern networks are known to be
miscalibrated without correction~\cite{guo2017calibration}. Shapley
values~\cite{shapley1953value,lundberg2017shap} provide an axiomatic credit
assignment that we adapt to attribute answers to retrieved documents---the
technical basis of original-source tracking. Self-RAG~\cite{asai2024selfrag}
adds self-reflective critique. Our framework integrates calibrated ensembles
and Shapley attribution natively rather than as post-hoc add-ons.

\section{Clinical QA Resources}
DDXPlus~\cite{tchango2022ddxplus} provides large-scale synthetic patients with
evidences and differential diagnoses; MedMCQA~\cite{pal2022medmcqa} and
MedQA~\cite{jin2021medqa} provide exam-style questions with explanations. We use
these as, respectively, the dynamics-training source, the evidence corpus, and
the held-out evaluation set.

%==============================================================
\chapter{Requirements and Planning}
%==============================================================

\section{Final Specifications and Requirements}
Table~\ref{tab:requirements} lists the functional and non-functional
requirements that the system must satisfy. The specifications follow directly
from the trustworthiness objectives: every answer must carry context labels and
source attribution, uncertainty must be calibrated, and the whole system must be
reproducible on commodity hardware.

\begin{table}[H]
\caption{Functional (F) and non-functional (N) requirements.}
\label{tab:requirements}
\centering
\small
\begin{tabular}{llp{8.2cm}}
\toprule
\textbf{ID} & \textbf{Type} & \textbf{Requirement}\\
\midrule
R1 & F & Hybrid (dense $+$ sparse) retrieval with fusion and re-ranking.\\
R2 & F & Context labels: evidence set, graph relations, predicted state,
         uncertainty.\\
R3 & F & Original-source tracking via per-document Shapley attribution and
         citations.\\
R4 & F & Calibrated uncertainty with epistemic/aleatoric decomposition.\\
R5 & F & Prospective planning via a latent world model and beam search.\\
R6 & N & Compact core ($\le$$15$M parameters); train in $<$$30$ min on one GPU.\\
R7 & N & Calibration: ECE $\le 0.05$; coverage near nominal.\\
R8 & N & Reproducibility: released code, configs, and figure scripts.\\
\bottomrule
\end{tabular}
\end{table}

\section{Societal Impact}
A clinical assistant that exposes calibrated confidence and traceable sources
can reduce automation bias by signalling when its output should be questioned,
and can widen access to structured clinical reasoning in settings with few
specialists. Conversely, any decision-support tool risks over-reliance.
Context labelling and source tracking are deliberate countermeasures: by making
the basis of each answer visible and auditable, the system supports
\emph{informed} use rather than blind trust, and keeps a human clinician in the
loop as the final decision-maker.

\section{Environmental Impact}
The neural core is intentionally small ($10.1$M parameters) and trains in under
twenty minutes on a single consumer GPU, so its training footprint is negligible
compared with large foundation-model training. At inference, the local core adds
only a few milliseconds per query; the dominant cost is the external LLM call,
which the framework invokes once per answer. The economy of the design
(Section~\ref{sec:economic}) is therefore also an environmental advantage:
fewer and smaller computations per trustworthy answer.

\section{Ethical Issues}
The work uses only public, de-identified datasets (DDXPlus, MedMCQA, MedQA) and
does not process real patient data. The system is a decision-\emph{support}
tool, not an autonomous diagnostician: it surfaces evidence, confidence, and
provenance so that a qualified clinician makes the final judgement. Because the
training trajectories are semi-synthetic and the outcomes are
differential-diagnosis probabilities, we are careful not to narrate them as
survival or treatment guarantees (Section~\ref{sec:discussion}). Calibration and
source tracking are themselves ethical safeguards against confident,
unattributable error.

\section{Standards}
The design follows established practice for the components it reuses: BM25 for
lexical retrieval, reciprocal rank fusion for list combination, cross-encoder
re-ranking, maximal marginal relevance for diversity, deep ensembles for
uncertainty, and Shapley values for attribution. Expected calibration error is
reported following the standard binning protocol~\cite{guo2017calibration}, and
statistical significance uses the Wilcoxon signed-rank test. Reproducibility
follows open-science norms: code, configurations, and scripts that regenerate
every figure and table are released.

\section{Project Management Plan}
The project was executed over an eighteen-week plan
(Fig.~\ref{fig:gantt}) spanning literature review, data collection and cleaning,
preprocessing and integration, model design and implementation, training and
tuning, evaluation and statistics, comparison and analysis, and writing.
Phases overlap where dependencies allow---for example, preprocessing began
before model implementation was complete.

\begin{figure}[H]
  \centering
  \includegraphics[width=0.92\textwidth]{fig21_gantt}
  \caption{Project management plan. Eight phases over an eighteen-week schedule,
  with deliberate overlap between data, modelling, and evaluation tasks.}
  \label{fig:gantt}
\end{figure}

\section{Risk Management}
Six principal risks were tracked on a likelihood--impact matrix
(Fig.~\ref{fig:risk}). The highest-impact risk (R4, overconfident
hallucination) is mitigated directly by the calibrated ensemble and the
hallucination check; data-related risks (R1 semi-synthetic trajectories, R2
class imbalance) are mitigated by honest framing and by reporting calibration
and coverage rather than relying on macro-F1; cost and compute risks (R3, R5)
are mitigated by the compact core and a single LLM call per query; and
reproducibility (R6) is mitigated by releasing all code and configurations.

\begin{figure}[H]
  \centering
  \includegraphics[width=0.95\textwidth]{fig22_risk_matrix}
  \caption{Risk assessment matrix. Markers R1--R6 are positioned by likelihood
  and impact; the side legend names each risk. Mitigations are described in the
  text.}
  \label{fig:risk}
\end{figure}

\section{Economic Analysis}
\label{sec:economic}
The direct project cost is dominated by external LLM API usage during
evaluation; GPU training is brief and the datasets are open
(Fig.~\ref{fig:economic}a). At the margin, the framework's local core lowers the
per-query cost relative to a cloud-LLM-only design because the small local model
absorbs the structured reasoning, leaving a single LLM call for verbalisation
(Fig.~\ref{fig:economic}b). The result is a favourable cost profile for a
trustworthy answer.

\begin{figure}[H]
  \centering
  \includegraphics[width=0.95\textwidth]{fig23_economic}
  \caption{Economic analysis. (a) Direct project cost by category. (b) Marginal
  per-query inference cost of the proposed local-core-plus-LLM design versus a
  cloud-LLM-only baseline.}
  \label{fig:economic}
\end{figure}

%==============================================================
\chapter{Design and Implementation}
%==============================================================

\section{Design Process and Methodology Overview}
The framework is organised into five functional blocks, mapped onto the detailed
architecture of Fig.~\ref{fig:archdetail}: hybrid graph-aware retrieval, a causal
knowledge-graph encoder, a latent world model, calibrated uncertainty, and
Shapley source tracking. This section specifies each block and the training
objectives.

\subsection{Hybrid Graph-Aware Retrieval}
Given a query $q$, HyDE~\cite{gao2023hyde} generates $n_h{=}3$ hypothetical
sub-answers that are embedded and averaged to form an expanded query
representation. We retrieve with a dense MiniLM bi-encoder over a
FAISS~\cite{johnson2019faiss} index and with BM25~\cite{robertson2009bm25}, and
fuse the two ranked lists by reciprocal rank fusion
(RRF)~\cite{cormack2009rrf},
\begin{equation}
\mathrm{RRF}(d)=\sum_{r\in\{\text{dense},\text{bm25}\}}\frac{1}{k+\mathrm{rank}_r(d)},
\quad k=60 .
\end{equation}
A cross-encoder~\cite{nogueira2019passage} re-scores the top candidates, and
maximal marginal relevance~\cite{carbonell1998mmr} selects a final set that
balances relevance and diversity,
\begin{equation}
\mathrm{MMR}=\arg\max_{d\in\mathcal{R}\setminus\mathcal{S}}
\big[\lambda\,\mathrm{rel}(d,q)-(1{-}\lambda)\max_{d'\in\mathcal{S}}\mathrm{sim}(d,d')\big],
\end{equation}
with $\lambda{=}0.6$. The selected evidence set $D$ becomes the first context
label and the candidate pool for source tracking.

\subsection{Causal Knowledge-Graph Encoder}
A medical knowledge graph with $N{=}20$ nodes encodes a sparse stochastic causal
prior over symptom--pathology relations. An edge-biased graph-attention
network~\cite{velickovic2018gat} processes the graph; for heads $h$ and edge
weights $e_{ij}$ the attention logits are
\begin{equation}
\alpha_{ij}^{(h)}\propto\frac{(\mathbf{W}_q x_i)^\top(\mathbf{W}_k x_j)}{\sqrt{d_h}}
+ \mathbf{W}_e\, e_{ij},
\end{equation}
with self-loops added for numerical stability. A masked mean pool yields a graph
embedding $\mathbf{g}\in\mathbb{R}^{256}$. The patient trajectory
$\mathbf{obs}\in\mathbb{R}^{T\times64}$ is encoded by an evidence-fusion
Transformer~\cite{vaswani2017attention} with gated cross-attention to
$\mathbf{g}$, producing latent states $\zt\in\mathbb{R}^{256}$ that serve as the
predicted-state context label.

\subsection{Latent World Model}
The core contribution is an action-conditioned latent dynamics model. At each
step a \emph{causal residual} computes the effect of action $a_t$ on the state,
\begin{equation}
\Delta\zt = s\cdot \sigma\!\big(\mathbf{W}_g[\zt,\mathbf{h}]\big)\odot
\mathrm{MLP}\big([\zt,\mathbf{g},\mathbf{e}_{a_t}]\big),
\end{equation}
where $\mathbf{e}_{a_t}$ is an action embedding and $s$ a learned scale. A
three-layer gated recurrent unit~\cite{cho2014gru} integrates the corrected
state and action to predict the next latent state,
\begin{equation}
\zt[t+1] = \mathrm{h2z}\big(\mathrm{GRU}([\zt+\Delta\zt,\ \mathbf{e}_{a_t}])\big).
\end{equation}
Three heads decode the latent: a \emph{decoder} reconstructs the observation, a
\emph{$\sigma$-head} emits per-dimension aleatoric uncertainty, and a
\emph{reward head} scores each transition. The model can roll out purely in
latent space, which is what enables prospective planning.

\subsection{Tree-of-Thought Beam Planner}
We treat plan selection as a search over the action tree with the world model as
transition function. From the encoded state $\mathbf{z}_0$, a beam search of
width $W{=}8$ and depth $H{=}4$ expands each beam node over the action space,
scores the resulting transition with the trained reward head, and retains the
top-$W$ partial plans. Terminal plans are ranked by a composite objective,
\begin{equation}
J(\pi) = v_\theta(\mathbf{z}_H) + \beta\!\sum_{t}r_\phi(\zt,\zt[t+1])
- \gamma\,\bar{\sigma}(\mathbf{z}_H),
\end{equation}
with $\beta{=}0.5$, $\gamma{=}0.1$, replacing undirected random rollouts with a
deliberate, reward-guided, uncertainty-aware search.

\subsection{Deep-Ensemble Uncertainty}
Five probabilistic heads~\cite{lakshminarayanan2017ensembles}, each emitting a
Gaussian $(\mu_m,\sigma_m^2)$ over the $5$-dimensional outcome (the
differential-diagnosis distribution), are combined as
\begin{equation}
\mu=\tfrac{1}{M}\sum_m\mu_m,\quad
\underbrace{\mathrm{Var}_m(\mu_m)}_{\text{epistemic}}+
\underbrace{\tfrac{1}{M}\sum_m\sigma_m^2}_{\text{aleatoric}}=\sigma^2_{\text{total}} .
\end{equation}
The decomposition separates reducible model uncertainty from irreducible data
noise and forms the calibrated-uncertainty context label.

\subsection{Shapley Source Tracking}
Original-source tracking attributes the answer to each retrieved document $d_i$
by its Shapley value~\cite{shapley1953value}, estimated by Monte-Carlo
permutation sampling~\cite{lundberg2017shap} with $32$ permutations,
\begin{equation}
\phi_i=\mathbb{E}_{\pi}\big[v(S_i^\pi\cup\{i\})-v(S_i^\pi)\big],
\end{equation}
where $v(\cdot)$ is a cross-encoder relevance score of the document subset to
$q$. The signed $\phi_i$ reveal supporting versus distracting evidence and are
surfaced as per-claim citations.

\subsection{Training Objectives}
The world model, GNN, and encoder are trained jointly to minimise a
reconstruction loss with latent-smoothness and uncertainty regularisers and a
supervised reward objective,
\begin{equation}
\mathcal{L}_{\text{wm}} = \underbrace{\|\hat{o}-o^+\|^2}_{\text{recon}}
+ 0.01\,\|\Delta\mathbf{z}\|^2
+ 0.001\,\bar{\sigma}
+ 0.1\,\|r_\phi-r^\star\|^2 .
\end{equation}
The reward target $r^\star_t = (\max_k y_k)\cdot t/T$ shapes the reward head to
prefer trajectories that drive the state toward a confident differential. The
ensemble is trained by Gaussian negative log-likelihood,
\begin{equation}
\mathcal{L}_{\text{ens}}=\tfrac{1}{M}\sum_m\tfrac{1}{2}
\Big(\log\sigma_m^2+\frac{(y-\mu_m)^2}{\sigma_m^2}\Big).
\end{equation}
Both stages use AdamW~\cite{kingma2015adam} with mixed-precision training.
Key hyperparameters are listed in Table~\ref{tab:hyper}.

\begin{table}[H]
\caption{Key hyperparameters.}
\label{tab:hyper}
\centering
\small
\begin{tabular}{lll}
\toprule
\textbf{Group} & \textbf{Parameter} & \textbf{Value}\\
\midrule
\multirow{4}{*}{Architecture}
  & latent / hidden dim         & $256$ / $512$ \\
  & Transformer layers / heads  & $3$ / $8$ \\
  & GNN layers (edge-biased GAT) & $3$ \\
  & ensemble members            & $5$ \\
\midrule
\multirow{3}{*}{Planning}
  & beam width $W$ / depth $H$   & $8$ / $4$ \\
  & reward weight $\beta$        & $0.5$ \\
  & uncertainty penalty $\gamma$ & $0.1$ \\
\midrule
\multirow{4}{*}{Retrieval}
  & top-$k$ / embed dim          & $6$ / $384$ \\
  & MMR $\lambda$ / RRF $k$       & $0.6$ / $60$ \\
  & HyDE sub-queries             & $3$ \\
  & SHAP permutations            & $32$ \\
\midrule
\multirow{3}{*}{Optimisation}
  & WM / ensemble epochs         & $15$ / $10$ \\
  & WM / ensemble LR             & $2\!\times\!10^{-4}$ / $10^{-3}$ \\
  & batch size / precision       & $64$ / FP16 \\
\bottomrule
\end{tabular}
\end{table}

\section{Data Collection}
We assemble $100{,}000$ samples from three public sources
(Table~\ref{tab:data}). DDXPlus~\cite{tchango2022ddxplus} supplies patient
records that drive world-model dynamics; MedMCQA~\cite{pal2022medmcqa} items
become the retrievable evidence corpus; and MedQA-USMLE~\cite{jin2021medqa}
provides held-out evaluation queries. All three are loaded from public mirrors
without credentials and cached locally.

\begin{table}[H]
\caption{Dataset composition ($100{,}000$ samples).}
\label{tab:data}
\centering
\small
\begin{tabular}{llr}
\toprule
\textbf{Dataset} & \textbf{Role} & \textbf{Size}\\
\midrule
DDXPlus (train/val/test) & World-model dynamics & $80\text{k}/10\text{k}/10\text{k}$\\
MedMCQA                  & Evidence corpus      & $50{,}000$ docs\\
MedQA-USMLE              & Held-out evaluation  & $1{,}000$ queries\\
\bottomrule
\end{tabular}
\end{table}

\subsection{Data Cleaning}
Each DDXPlus record is parsed and validated before use. We audit the training
split for missing pathology labels, empty differential diagnoses, empty evidence
lists, and parse failures of the string-encoded list fields. All $80{,}000$
training records pass these completeness checks (Fig.~\ref{fig:cleaning}a),
confirming that the curated source requires correction of \emph{format} rather
than removal of \emph{records}. The exploratory analysis
(Fig.~\ref{fig:overview}) shows the cleaned distributions: patient age is
broadly spread (mean $39.6$, standard deviation $22.7$ years), the sex split is
near-balanced ($48.6\%$ male, $51.4\%$ female), the pathology distribution is
led by upper-respiratory and viral conditions, and patients carry a median of
$19$ clinical evidences each.

\begin{figure}[H]
  \centering
  \includegraphics[width=\textwidth]{fig13_data_overview}
  \caption{Exploratory data analysis of the cleaned DDXPlus training split:
  (a) patient age distribution; (b) sex distribution; (c) the fifteen most
  frequent pathologies; and (d) the number of clinical evidences per patient.}
  \label{fig:overview}
\end{figure}

\begin{figure}[H]
  \centering
  \includegraphics[width=\textwidth]{fig14_data_cleaning}
  \caption{Data cleaning. (a) Record-level validation funnel: all $80{,}000$
  training records satisfy the completeness checks. (b) Evidence-code
  normalisation collapses $504$ raw, value-suffixed codes to $223$ canonical
  base codes ($-56\%$), reducing sparsity before encoding.}
  \label{fig:cleaning}
\end{figure}

\subsection{Data Transformation}
Cleaned records are transformed into model-ready form. Raw evidence codes of the
form \texttt{E\_54\_@\_V\_161} are normalised to canonical base codes
(\texttt{E\_54}), which reduces the unique-code vocabulary from $504$ to $223$
(Fig.~\ref{fig:cleaning}b) and removes value-suffix sparsity. Age is scaled to
$[0,1]$, sex is mapped to a binary indicator, and the evidence set is encoded as
binary flags. Each patient is expanded into a $T{=}8$-step diagnostic trajectory
in which evidence is revealed sequentially, yielding a $64$-dimensional
observation per step, a discrete evidence-acquisition action, and a
$5$-dimensional outcome equal to the top-$k$ differential-diagnosis
distribution.

\subsection{Data Integration}
The three sources are integrated into a single pipeline: DDXPlus trajectories
train and evaluate the world model and ensemble; MedMCQA items are embedded and
indexed (dense $+$ BM25) as the retrievable evidence corpus that supplies
context labels and the Shapley attribution pool; and MedQA-USMLE queries drive
the end-to-end comparison. A shared vocabulary of $49$ pathologies and $223$
evidence base codes links the structured and textual modalities.

\subsection{Data Reduction}
The preprocessing pipeline (Fig.~\ref{fig:preprocess}) reduces high-cardinality,
string-encoded records to compact tensors: cleaning removes format noise,
transformation normalises and encodes, integration joins the sources, and
reduction projects each patient onto a fixed $64$-dimensional observation space
with a $5$-dimensional outcome. The result is a dense numerical representation
small enough to train in minutes yet rich enough to preserve clinically
meaningful structure.

\begin{figure}[H]
  \centering
  \includegraphics[width=\textwidth]{fig18_preprocess_pipeline}
  \caption{Data preprocessing pipeline: cleaning, transformation, integration,
  and reduction map raw clinical records to fixed-size tensors.}
  \label{fig:preprocess}
\end{figure}

\subsection{Summary of Preprocessed Data}
The preprocessed dataset exhibits two properties that shape evaluation. First,
the pathology distribution is strongly long-tailed across $49$ classes
(Fig.~\ref{fig:imbalance}a), and the true pathology is the top-ranked entry of
the differential for the large majority of patients
(Fig.~\ref{fig:imbalance}b); this explains why rank-based macro-F1 is a
demanding metric. Second, the engineered $64$-dimensional observation space is
informative: the age and sex dimensions are almost always active, evidence
dimensions activate at clinically plausible rates, and the evidence
co-occurrence structure forms coherent blocks (Fig.~\ref{fig:preprocessed}),
indicating that the encoding preserves real symptom correlations rather than
noise.

\begin{figure}[H]
  \centering
  \includegraphics[width=\textwidth]{fig15_class_imbalance}
  \caption{Class structure of the preprocessed data. (a) Pathology frequency
  across $49$ classes on a log scale, showing a long tail. (b) Distribution of
  the true pathology's rank within the differential-diagnosis target.}
  \label{fig:imbalance}
\end{figure}

\begin{figure}[H]
  \centering
  \includegraphics[width=\textwidth]{fig16_preprocessed_summary}
  \caption{Preprocessed feature space. (a) Activation frequency of each of the
  $64$ observation dimensions (age and sex highlighted). (b) Pearson
  co-occurrence among the eighteen most active evidence dimensions, showing
  coherent clinical blocks.}
  \label{fig:preprocessed}
\end{figure}

%==============================================================
\chapter{Result Analysis}
%==============================================================

\section{Performance Evaluation}
\label{sec:performance}
We evaluate the world model and ensemble on the DDXPlus validation split using
next-state mean absolute error (MAE) and RMSE, $1\sigma$ coverage, expected
calibration error (ECE)~\cite{guo2017calibration}, top-rank diagnosis accuracy,
and macro-F1. Table~\ref{tab:main} and Fig.~\ref{fig:metrics} summarise the
results. The model predicts the next latent-decoded state with MAE $0.039$ and
RMSE $0.074$, ranks the true pathology at the top of the differential with
$0.755$ accuracy, and is well calibrated, with ECE $0.029$ and $1\sigma$
coverage $0.80$. The macro-F1 of $0.17$ reflects the difficulty and class
imbalance of fine-grained rank classification across $49$ pathologies
(Section~\ref{sec:discussion}).

\begin{table}[H]
\caption{World-model and ensemble evaluation on the DDXPlus validation split.}
\label{tab:main}
\centering
\small
\begin{tabular}{lc}
\toprule
\textbf{Metric} & \textbf{Value}\\
\midrule
Next-state MAE        & $0.039$ \\
Next-state RMSE       & $0.074$ \\
Top-rank accuracy     & $0.755$ \\
Macro-F1              & $0.172$ \\
Expected calibration error (ECE) & $0.029$ \\
$1\sigma$ coverage    & $0.801$ \\
Mean $|\text{SHAP}|$  & $0.734$ \\
\midrule
Parameters            & $10{,}130{,}060$ \\
Training time         & $<20$ min (1$\times$RTX 4080S) \\
\bottomrule
\end{tabular}
\end{table}

\begin{figure}[H]
  \centering
  \includegraphics[width=0.95\textwidth]{fig7_metrics_dashboard}
  \caption{World-model evaluation metrics. Low error and calibration error
  accompany competitive accuracy and well-behaved coverage.}
  \label{fig:metrics}
\end{figure}

\paragraph{Generalisation (testing method).}
To verify that the model does not overfit, we evaluate the trained checkpoints on
a held-out split and compare against training-split performance
(Fig.~\ref{fig:traintest}). Training and held-out errors are nearly identical
(MAE $0.037$ vs.\ $0.039$; RMSE $0.070$ vs.\ $0.074$; error rate $0.232$ vs.\
$0.245$), a gap below $0.002$ MAE, confirming strong generalisation. The
training curves (Fig.~\ref{fig:traintest}a) show smooth convergence of the joint
loss, the reconstruction term, and the newly supervised reward head.

\begin{figure}[H]
  \centering
  \includegraphics[width=\textwidth]{fig19_train_test}
  \caption{Training and generalisation. (a) World-model training convergence
  (total, reconstruction, and reward-head losses). (b) Train versus held-out
  error: the small gap indicates the model generalises rather than memorises.}
  \label{fig:traintest}
\end{figure}

\paragraph{Calibration and uncertainty (reliability).}
Figure~\ref{fig:calibration} shows the reliability diagram: binned predicted
uncertainty tracks empirical error along the diagonal, confirming that the
ensemble's confidence is meaningful. The epistemic/aleatoric decomposition
(Fig.~\ref{fig:uncertainty}) separates reducible model uncertainty from
irreducible data noise---information the planner exploits through the
$\gamma$-penalty and which the system reports as a context label.

\begin{figure}[H]
  \centering
  \includegraphics[width=0.7\textwidth]{fig4_calibration}
  \caption{Reliability diagram. Predicted uncertainty closely matches empirical
  error, indicating good calibration (low ECE).}
  \label{fig:calibration}
\end{figure}

\begin{figure}[H]
  \centering
  \includegraphics[width=0.74\textwidth]{fig5_uncertainty}
  \caption{Epistemic versus aleatoric uncertainty over the validation set. The
  spread shows the ensemble distinguishes model uncertainty from data noise.}
  \label{fig:uncertainty}
\end{figure}

\paragraph{Latent structure.}
The t-SNE projection of the learned latent states (Fig.~\ref{fig:latent})
exhibits organisation by true diagnosis rank, indicating that the encoder and
dynamics model capture clinically meaningful geometry rather than collapsing to a
trivial embedding.

\begin{figure}[H]
  \centering
  \includegraphics[width=0.72\textwidth]{fig6_latent_tsne}
  \caption{t-SNE of latent states, coloured by the true diagnosis rank. The
  emergent organisation reflects learned clinical structure.}
  \label{fig:latent}
\end{figure}

\section{Analysis of Design Solutions}
The design couples a mature retrieval stack with a compact world-model core. Its
strengths are calibration (ECE $0.029$), generalisation (train/test gap $<0.002$
MAE), efficiency ($10.1$M parameters, $<20$ min training), and native
trustworthiness (context labels and source tracking in a single pass). Its
principal weaknesses follow from the data: semi-synthetic trajectories make the
transition target close to deterministic, and severe class imbalance depresses
macro-F1. The feasibility and cost analysis (Sections~\ref{sec:economic}
and~\ref{sec:performance}) favours the design: the small local core is cheap to
train and run, and it lowers the marginal per-query cost relative to a
cloud-LLM-only system while adding the trust mechanisms the baseline lacks.

\section{Final Design Adjustments}
Two adjustments were made in response to the performance evaluation. First, the
reward head was changed from unsupervised to \emph{supervised} with the shaping
target $r^\star_t=(\max_k y_k)\cdot t/T$, so that planning scores became
meaningful rather than arbitrary; the reward loss then converges alongside
reconstruction (Fig.~\ref{fig:traintest}a). Second, the planner's undirected
random rollouts were replaced by a \emph{beam search} ($W{=}8$, $H{=}4$) that
maximises the composite objective $J(\pi)$. Neither change alters the calibrated
predictive metrics, which depend on the ensemble, but together they make the
planning objective well defined.

\section{Statistical Analysis}
\label{sec:stat}
To test whether the improvement over the baseline is significant rather than
incidental, we apply the Wilcoxon signed-rank test to the paired per-vignette
concept-coverage scores of the two systems across the ten clinical vignettes
(Fig.~\ref{fig:stat}). The proposed framework improves coverage on every paired
case (Fig.~\ref{fig:stat}a); the mean improvement is $+0.27$ with a $95\%$
confidence interval of $\pm0.21$, and the signed-rank test rejects the null
hypothesis of no difference at $p=0.046$ (Fig.~\ref{fig:stat}b). The effect is
therefore statistically significant at the $5\%$ level despite the modest sample
of vignettes.

\begin{figure}[H]
  \centering
  \includegraphics[width=\textwidth]{fig20_statistical}
  \caption{Statistical analysis of the controlled comparison. (a) Paired
  per-vignette concept-coverage scores (baseline vs.\ proposed). (b) Mean
  improvement with a $95\%$ confidence interval and the Wilcoxon signed-rank
  $p$-value.}
  \label{fig:stat}
\end{figure}

\section{Comparisons and Relationships}
We compare the framework against representative RAG families and against a
strong measured baseline. Qualitatively (Fig.~\ref{fig:capability}), existing
methods---Vanilla RAG~\cite{lewis2020rag}, HyDE~\cite{gao2023hyde},
cross-encoder re-ranking RAG~\cite{nogueira2019passage}, and
Self-RAG~\cite{asai2024selfrag}---cover retrieval and re-ranking but lack causal
reasoning, world-model planning, calibrated uncertainty, per-document
attribution, and treatment simulation, all of which the proposed framework
provides in a single pipeline.

Quantitatively, the control is a strong hybrid RAG system (dense MiniLM$+$FAISS
and BM25 fused by RRF), answered directly by the same DeepSeek
LLM~\cite{deepseek2024} used by our framework, but without the world model,
ensemble, or attribution. Both systems answer an identical set of ten complex
clinical vignettes over the same corpus. Table~\ref{tab:compare},
Fig.~\ref{fig:comparison}, and Fig.~\ref{fig:quantcompare} summarise the study.
The framework raises clinical concept coverage from $0.70$ to $0.97$ ($+38\%$),
answer-structure completeness from $0.50$ to $0.84$ ($+68\%$), and evidence
grounding from $2.0$ to $5.1$ citations per answer ($2.5\times$). Per-vignette
results (Fig.~\ref{fig:perprompt}) show the gain is consistent across all ten
cases. Notably, the framework reports \emph{lower but better-calibrated}
confidence ($0.70$ vs.\ $0.91$): the baseline's near-constant high confidence is
characteristic of overconfidence, whereas our framework tempers confidence in
proportion to genuine uncertainty. Finally, the framework supplies a mean
$|\text{SHAP}|$ attribution of $1.28$ for source tracking, calibrated
uncertainty, and treatment-plan simulation---none of which the baseline
provides.

\begin{table}[H]
\caption{Proposed framework versus hybrid-RAG baseline on ten clinical
vignettes.}
\label{tab:compare}
\centering
\small
\begin{tabular}{lcc}
\toprule
\textbf{Metric} & \textbf{Baseline RAG} & \textbf{Proposed}\\
\midrule
Clinical concept coverage      & $0.70$ & $\mathbf{0.97}$\\
Answer-structure completeness  & $0.50$ & $\mathbf{0.84}$\\
Evidence citations (avg.)      & $2.0$  & $\mathbf{5.1}$\\
Stated confidence              & $0.91$ & $0.70$ (calibrated)\\
SHAP attribution (mean $|\phi|$)& ---   & $\mathbf{1.28}$\\
Calibrated uncertainty         & No     & \textbf{Yes}\\
Plan simulation                & No     & \textbf{Yes}\\
\bottomrule
\end{tabular}
\end{table}

\begin{figure}[H]
  \centering
  \includegraphics[width=\textwidth]{fig11_method_capability}
  \caption{Capability comparison against representative RAG families. Existing
  methods cover retrieval and re-ranking but lack causal reasoning, world-model
  planning, calibrated uncertainty, per-document attribution, and treatment
  simulation, all of which the proposed framework provides.}
  \label{fig:capability}
\end{figure}

\begin{figure}[H]
  \centering
  \includegraphics[width=0.92\textwidth]{fig8_comparison}
  \caption{Aggregate comparison on ten vignettes. The proposed framework
  improves coverage, completeness, and grounding over the baseline; citation
  counts are shown normalised by the top-$k$ budget.}
  \label{fig:comparison}
\end{figure}

\begin{figure}[H]
  \centering
  \includegraphics[width=0.92\textwidth]{fig12_method_quantitative}
  \caption{Head-to-head quantitative comparison against the strongest measured
  baseline (hybrid dense$+$BM25 RAG) on the ten vignettes. The proposed
  framework improves concept coverage, structure completeness, and evidence
  grounding, and additionally provides calibrated confidence---a capability the
  baseline lacks (shown as available$=1$/unavailable$=0$).}
  \label{fig:quantcompare}
\end{figure}

\begin{figure}[H]
  \centering
  \includegraphics[width=0.92\textwidth]{fig9_perprompt}
  \caption{Per-vignette clinical concept coverage. The proposed framework
  matches or exceeds the baseline on every case.}
  \label{fig:perprompt}
\end{figure}

\section{Multi-Method Comparison Against Advanced RAG}
\label{sec:multirag}
The head-to-head study above isolates the contribution of the world-model core
against a single strong baseline. To situate the framework within the broader
landscape of \emph{advanced} retrieval, we additionally run four progressively
stronger RAG pipelines \emph{live} through the same DeepSeek
LLM~\cite{deepseek2024}, over the same MedMCQA evidence corpus, on the same ten
clinical vignettes, so that the only variable is the retrieval/reasoning stack:
(M1)~vanilla dense RAG (MiniLM$+$FAISS); (M2)~hybrid dense$+$BM25 retrieval with
reciprocal-rank fusion; (M3)~HyDE~\cite{gao2023hyde} hypothetical-document
expansion over hybrid retrieval; and (M4)~cross-encoder
re-ranking~\cite{nogueira2019passage} with maximal-marginal-relevance (MMR)
diversification. Method~M5 is the full proposed framework.
Table~\ref{tab:multirag} and Fig.~\ref{fig:multirag} report the measured study.

\begin{table}[H]
\caption{Live multi-method comparison of four advanced-RAG pipelines and the
proposed framework on ten clinical vignettes (identical corpus, LLM, and scoring
rubric; only the retrieval/reasoning stack varies). Higher is better except
stated confidence, which the framework deliberately calibrates rather than
maximises.}
\label{tab:multirag}
\centering
\small
\begin{tabular}{lcccc}
\toprule
\textbf{Method} & \textbf{Concept} & \textbf{Structure} & \textbf{Citations} & \textbf{Conf.}\\
\midrule
M1\quad Vanilla dense RAG          & $0.80$ & $0.51$ & $2.3$ & $0.86$\\
M2\quad Hybrid (dense$+$BM25)      & $0.73$ & $0.53$ & $1.6$ & $0.84$\\
M3\quad HyDE $+$ hybrid            & $0.87$ & $0.50$ & $2.5$ & $0.83$\\
M4\quad Cross-encoder $+$ MMR      & $0.77$ & $0.51$ & $1.5$ & $0.91$\\
\midrule
\textbf{M5\quad Proposed (ours)}   & $\mathbf{0.97}$ & $\mathbf{0.84}$ & $\mathbf{5.1}$ & $0.70$\\
\bottomrule
\end{tabular}
\end{table}

\begin{figure}[H]
  \centering
  \includegraphics[width=\textwidth]{fig24_multi_method}
  \caption{Advanced-RAG methods versus the proposed framework on the ten
  vignettes (identical corpus and LLM). Left: answer quality (concept coverage,
  structure completeness, and evidence grounding normalised by the top-$k$
  budget). Right: end-to-end latency---HyDE and cross-encoder$+$MMR add an extra
  LLM call and a re-ranking pass without closing the quality gap.}
  \label{fig:multirag}
\end{figure}

Two findings stand out. First, stronger \emph{retrieval} alone yields only
incremental and non-monotonic gains: HyDE attains the best retrieval-only
concept coverage ($0.87$), yet hybrid fusion and MMR re-ranking do not uniformly
help on these complex vignettes---hybrid even trails vanilla dense retrieval on
coverage ($0.73$ vs.\ $0.80$)---and every retrieval-only variant plateaus near
$0.50$ on answer-structure completeness. Second, the largest and most
\emph{trust-relevant} improvements come from the reasoning stack layered on top
of retrieval: the framework lifts concept coverage to $0.97$ ($+10$ points over
the strongest baseline, HyDE), nearly doubles structural completeness to $0.84$,
and roughly doubles evidence grounding to $5.1$ citations per answer. The
advanced retrieval methods also incur higher latency (HyDE and
cross-encoder$+$MMR add $4$--$5$\,s per query for the extra LLM call and
re-ranking pass; Fig.~\ref{fig:multirag}, right) without closing the quality
gap. Crucially, none of M1--M4 provides calibrated uncertainty, per-document
SHAP attribution, or treatment-plan simulation, all of which the framework
supplies in the same pass.

\section{Worked Answer Comparisons}
\label{sec:demo}
To make the quantitative gains concrete, we show the actual model outputs on
three representative vignettes, with the measured per-answer scores. Excerpts are
trimmed for length (``\ldots'') but otherwise verbatim.

\paragraph{Vignette P1 --- Acute coronary syndrome.}
\emph{58-year-old male with hypertension and diabetes; acute chest pain
radiating to the left arm and jaw; ST-elevation in leads II, III, aVF; troponin
$2.5$\,ng/mL; pain for $45$ minutes.}

\noindent
\begin{minipage}[t]{0.47\textwidth}
\footnotesize
\textbf{Baseline RAG}\\{\scriptsize coverage $0.67$, $2$ citations, conf.\ $0.95$}\\[3pt]
``Most likely diagnosis: inferior-wall STEMI. Immediate management: emergency
reperfusion therapy (primary PCI or thrombolysis) [2][3].''
\end{minipage}\hfill
\begin{minipage}[t]{0.47\textwidth}
\footnotesize
\textbf{GRAPES-SHAP (ours)}\\{\scriptsize coverage $1.00$, $4$ citations, conf.\ $0.70$}\\[3pt]
Structured reasoning: confirms inferior-wall STEMI [2]; prioritises reperfusion
and explicitly weighs \textbf{PCI vs.\ thrombolysis} (PCI preferred --- lower
intracranial-haemorrhage risk); treatment choice informed by plan simulation
(score $-0.29$); adds the supportive-care bundle (aspirin, P2Y12 inhibitor,
heparin). Reports calibrated confidence $0.70$ and mean$|\text{SHAP}|=1.08$.
\end{minipage}

\smallskip
\noindent\textit{Takeaway:} both reach the correct top-line diagnosis, but only
the framework supplies the PCI-versus-thrombolysis rationale, the full
supportive-care bundle, calibrated confidence, and per-evidence attribution.

\paragraph{Vignette P2 --- Septic shock.}
\emph{72-year-old with COPD; fever, altered mental status, HR $118$, lactate
$3.2$\,mmol/L, BP $88/52$, pneumonia on imaging, cultures pending.}

\noindent
\begin{minipage}[t]{0.47\textwidth}
\footnotesize
\textbf{Baseline RAG}\\{\scriptsize coverage $0.33$, $1$ citation, conf.\ $0.90$}\\[3pt]
``Septic shock secondary to pneumonia; aggressive fluid resuscitation, oxygen,
and norepinephrine for hypotension if fluids fail [3].''
\end{minipage}\hfill
\begin{minipage}[t]{0.47\textwidth}
\footnotesize
\textbf{GRAPES-SHAP (ours)}\\{\scriptsize coverage $1.00$, $6$ citations, conf.\ $0.70$}\\[3pt]
Structured plan around the three pillars --- source control, \textbf{early
broad-spectrum antibiotics} after cultures, and haemodynamic support;
reconciles the retrieved ``dopamine-first'' text against modern
norepinephrine-first practice; grounds each step ($6$ citations).
\end{minipage}

\smallskip
\noindent\textit{Takeaway:} the baseline omits early broad-spectrum antibiotics
--- arguably the single most outcome-critical step --- which the framework
recovers while flagging the evidence-versus-guideline discrepancy.

\paragraph{Vignette P4 --- Acute ischaemic stroke.}
\emph{65-year-old with atrial fibrillation (not anticoagulated); acute aphasia
and right-sided weakness; last known normal $1.5$\,h ago; NIHSS $14$; CT shows
no haemorrhage.}

\noindent
\begin{minipage}[t]{0.47\textwidth}
\footnotesize
\textbf{Baseline RAG}\\{\scriptsize coverage $0.33$, $2$ citations, conf.\ $0.95$}\\[3pt]
``Acute ischaemic stroke (left MCA); administer IV thrombolysis (alteplase)
within $4.5$\,h of onset [2][3].''
\end{minipage}\hfill
\begin{minipage}[t]{0.47\textwidth}
\footnotesize
\textbf{GRAPES-SHAP (ours)}\\{\scriptsize coverage $1.00$, $4$ citations, conf.\ $0.70$}\\[3pt]
Traces localisation $\rightarrow$ cardioembolic AF aetiology $\rightarrow$
eligibility (within window, no contraindication), then selects IV thrombolysis
(alteplase/tenecteplase); mean$|\text{SHAP}|=1.36$ ranks the time-window and
no-haemorrhage evidence as decisive.
\end{minipage}

\smallskip
\noindent\textit{Takeaway:} both reach tPA, but only the framework exposes the
localisation$\rightarrow$aetiology$\rightarrow$eligibility chain and the
evidence weights behind the time-critical decision.

\section{Discussions}
\label{sec:discussion}
The results show that inserting an explicit, calibrated, and interpretable world
model between retrieval and generation---and surfacing its state as context
labels and source attributions---yields measurable, statistically significant
gains in grounding and answer quality at a small parameter and compute budget.
We nonetheless delineate the method's scope precisely, as responsible
clinical-AI reporting demands.

\textbf{Semi-synthetic trajectories.} DDXPlus is a static diagnostic dataset;
the preprocessor constructs trajectories by revealing evidence sequentially. The
transition target is therefore close to deterministic, and the low MAE should be
read as faithful modelling of this constructed dynamics, not of real
physiological time series.

\textbf{Action semantics.} In training, actions correspond to \emph{evidence
acquisition} rather than physical interventions. The planner is thus best
understood as an evidence-gathering / diagnostic-refinement policy; mapping it
onto treatment actions would require a dataset with logged interventions and
outcomes.

\textbf{Outcome interpretation.} The supervised outcomes are
differential-diagnosis probabilities. Narrating them as survival or readmission
would over-claim; we therefore report them as diagnostic confidence.

\textbf{Knowledge graph.} The graph encodes a stochastic causal \emph{prior}
rather than a curated ontology; substituting a literature-derived graph (e.g.,
SNOMED-CT relations) is a natural extension.

\textbf{Macro-F1.} The modest macro-F1 stems from fine-grained, imbalanced rank
classification (Fig.~\ref{fig:imbalance}) and the use of $\arg\max$ over a soft
distribution; it does not contradict the strong calibration and coverage
results, which are the quantities of interest for trustworthy decision support.

%==============================================================
\chapter{Conclusion}
%==============================================================

\section{Summary of Findings}
This thesis presented a framework that makes clinical question answering more
trustworthy through context labelling and original-source tracking, supported by
an action-conditioned latent world model, a calibrated deep ensemble, and
Shapley-value attribution in a compact $10.1$M-parameter model. The world model
is accurate (MAE $0.039$), well calibrated (ECE $0.029$), and generalises (train/test
gap $<0.002$ MAE). The full system significantly improves concept coverage,
completeness, and evidence grounding over a strong RAG baseline (Wilcoxon
$p=0.046$) while adding calibrated uncertainty and per-document source tracking
that the baseline lacks.

\section{Contributions to the Field}
The contributions are: (i) a trust-oriented reframing of RAG around context
labelling and original-source tracking; (ii) a compact world-model core with
supervised reward shaping and beam-search planning that adds prospective
reasoning to retrieval; (iii) native calibrated uncertainty with an
epistemic/aleatoric decomposition; (iv) Shapley-value source tracking that
attributes answers to individual documents; and (v) an honest, reproducible
characterisation of scope, with all code, configurations, and figure scripts
released.

\section{Recommendations for Future Work}
\label{sec:future}
The limitations define concrete next steps. Replacing the semi-synthetic
dynamics with logged intervention--outcome data (e.g., intensive-care records)
would enable true treatment planning. Substituting the stochastic causal prior
with a curated medical ontology (e.g., SNOMED-CT) would strengthen the graph
reasoning. Addressing class imbalance with re-weighting or focal objectives
would raise macro-F1 without sacrificing calibration. Finally, a prospective
clinician-in-the-loop study would test whether context labels and source
tracking measurably reduce automation bias in practice.

%==============================================================
% REPRODUCIBILITY
%==============================================================
\chapter*{Reproducibility}
\addcontentsline{toc}{chapter}{Reproducibility}
All code, configurations, trained-model metrics, and the scripts that generate
every figure and table in this thesis are released with the project. The
end-to-end retraining and evaluation are reproduced by
\texttt{scripts/retrain\_and\_eval.py}; the data statistics, exploratory
figures, and train/test analysis by \texttt{scripts/make\_paper\_assets.py}.
The comparison against the baseline reuses a single recorded LLM run and does not
require re-querying the external API.

%==============================================================
\bibliographystyle{IEEEtran}
\bibliography{references}

\end{document}
